\documentclass[10pt]{article}
\usepackage[margin=.8in]{geometry}
\usepackage[T1]{fontenc}
\usepackage{amsmath,amssymb,booktabs,longtable,hyperref,xurl}
\hypersetup{colorlinks=true,urlcolor=blue,linkcolor=blue}
\title{Potential-Initialized PCBF--CLF RTI: One Model per Initialization}
\author{Collision Avoidance Research}
\date{October 3, 2026}
\begin{document}
\maketitle
\section{Implemented behavior}
Normal frames shift the previous input sequence and roll it out from the current
state. They linearize that trajectory once and attempt the inherited-safety-budget
CLF problem. Startup, an unevaluable shift, or a failed shifted optimization uses
a fresh motion-aware potential-field initialization. A failed shifted model may
be replaced once by this new seed and its model. An optimized rollout is never
relinearized again within the same hold. A fresh seed that fails reports failure;
no previous plan, seed control or primary-only point becomes an executable backup.

The same single analytic quadratic CLF remains in every case. A zero PCBF lower
bound does not skip CLF optimization. Damping interpolates a feasible segment
of the same conic problem and uses cheap nonlinear model-agreement evaluations;
it does not rebuild dynamics or call another solver. Initialization/restoration
may still require PCBF and CLF solves. The numerical correction box can grow
between frames when agreement is good; permanently capping it at the small
initial scale substantially delayed nominal recovery in a discarded prototype.

\section{Transfer from Zhai et al.}
The local paper by Zhai, Liu, Zhang and Wang, IEEE Access 12 (2024),
19665--19681, DOI \path{10.1109/ACCESS.2024.3355952}, motivates directional repulsion,
relative-motion awareness, and coupled course/speed guidance (Sections III--IV).
This implementation does not reproduce its printed equations, tracked-vehicle
model, eight behavior priorities or Bezier tracker. The preceding assessment
\path{report/FLOW_INITIALIZATION_ZHAI_20261002.tex} records the formula issues.

At each seed hold, preview the actual ego tangent speed and an exponential
lateral return to the given path. Pair every ego preview position with the known
target at the same absolute time. The target keeps constant tangential
acceleration and sideslip, with exact constant-curvature propagation. In the
road frame, form the directional interaction measure
\[
 w_j=\exp\!\left[-\tfrac12\left((\Delta s_j/a_j)^2+
 (\Delta d_j/b_j)^2\right)-\tau_j/(2T)\right],\qquad w=\max_jw_j.
\]
The dimensions $a_j,b_j$ use rectangle supports, target orientation, body-center
offsets and seed padding. Lateral attraction, normalized lateral repulsion and
bounded circulation give a course demand; predicted longitudinal closing gives
a speed reduction. This is a guidance field, not the full gradient of a claimed
scalar safety potential. It never divides by closing speed. A passing sign is
held within each seed rather than enumerating maneuver modes.

Inverse Fiala front-force guidance and a speed loop produce controls; the actual
nonlinear bicycle RK4 map produces the anchor states. The final two seconds
settle intrinsic velocity/input state toward the existing free-pose straight
core. Endpoint position and heading remain free. The same completion is used
for target-free initialization, where regression testing exposed that continued
path recovery could miss the tiny core. No extra CLF or collision acceptance
test is added.

The seed parameters are common to all cases: three-second preview, 0.1-second
preview spacing, 0.7-m guidance padding, 0.3 braking-ratio shaping, 0.98 lateral
acceleration/front-force fractions, and unit speed gain. These are initialization
parameters, not physical steering limits or optimized-input clipping. The
braking parameter was selected through a bounded raw-seed sweep on the failed
accelerating head-on case. No claim of parameter optimality or generalization
outside the evaluated family is made. The old fitted Gaussian displacement,
transported-cylinder helper and obsolete comparison driver were removed.

\section{Experiment contract}
The seven prescribed collision threats are evaluated at 8 and 15 m/s. Holds are
50 ms; PCBF prefixes have 8/16 steps and the unchanged completion adds 60 steps.
The input is applied to ODE45 with relative/absolute tolerances $10^{-11}$ and
$10^{-12}$; 31 output samples per hold audit rectangle clearance. State is exact,
road constraints are absent, and target motion is fixed and known. There is no
random seed. Computation latency is measured but not injected into plant time.

The observation cap is 1,500 holds, not a success deadline. Success requires
post-conflict errors below
\[ [0.1\,\mathrm m,1^\circ,0.1\,\mathrm{m/s},
0.05\,\mathrm{m/s},0.01\,\mathrm{rad/s}] \]
continuously for one second.
The controller's 0.10-m margin constrains its affine collision rows. Dense
nonlinear replay clearance is a separate measurement, not an assertion that
this margin is certified on the continuous plant.

\section{Closed-loop results}
All 14 cases recovered without sampled collision or solver interruption, across 4,618 holds. Minimum replay clearance was 49.477 mm. The slowest recovery confirmation was 66.00 s.

\begin{center}\small
\begin{tabular}{rlrrrr}
\toprule
Speed & Scenario & Recovery (s) & Gap (m) & Max $|e_y|$ (m) & Run max (ms)\\
\midrule
8 & headOn & 12.60 & 1.7211 & 5.956 & 81.40\\
8 & acceleratingHeadOn & 12.55 & 1.3235 & 5.876 & 46.77\\
8 & brakingLead & 13.50 & 1.1830 & 5.549 & 52.86\\
8 & crossing & 12.60 & 0.0915 & 5.722 & 46.49\\
8 & turningCrossing & 12.10 & 0.1032 & 5.609 & 50.84\\
8 & curvedHeadOn & 17.90 & 2.1481 & 6.084 & 43.03\\
8 & curvedCrossing & 13.05 & 0.0944 & 5.750 & 51.53\\
15 & headOn & 10.15 & 0.1759 & 7.241 & 47.38\\
15 & acceleratingHeadOn & 10.30 & 0.0495 & 6.872 & 56.30\\
15 & brakingLead & 16.10 & 0.8721 & 7.871 & 56.21\\
15 & crossing & 11.55 & 3.3413 & 10.773 & 43.39\\
15 & turningCrossing & 11.40 & 3.1841 & 10.478 & 44.17\\
15 & curvedHeadOn & 66.00 & 3.9997 & 11.143 & 62.99\\
15 & curvedCrossing & 11.10 & 2.5831 & 10.017 & 44.56\\
\bottomrule
\end{tabular}\end{center}
``Run max'' excludes each scenario's first call. These are controller-call wall
measurements without a profiler; the workstation was shared with other research
processes. They establish an observed distribution, not a worst-case deadline.

Continuation median, 95th percentile and maximum were 13.572, 26.399 and 81.396 ms. There were 0 continuation frames over 100 ms and 18 over the 50-ms hold. The largest first call was 864.633 ms, including startup/cached-construction effects.

There were 4,545 one-model frames and 73 failed-shift/fresh-seed frames. 4,510 frames used one numerical solve, 33 used damping, and no frame used more than 3 solves. These counts distinguish one linearization per initialization from a promise of one numerical solve in every exception.

The slowest continuation was 8 m/s headOn, frame 5. Its measured components were:
\begin{center}\begin{tabular}{lr}\toprule Component & ms\\\midrule
Shift / initialization & 1.854\\
Problem construction (including CLF) & 32.987\\
Numerical solvers & 15.270\\
Full-step model agreement & 0.914\\
Damping & 0.000\\
Other controller overhead & 30.371\\
\bottomrule\end{tabular}\end{center}
Other overhead is a remainder, not a profiled function attribution.

The 66-second confirmation in the 15-m/s curved head-on case is not a
66-second first-encounter recovery. The trace first satisfies a one-second
nominal dwell at 11.1 s. The prescribed constant-sideslip target follows a
circular orbit and produces another encounter near 56 s (the cruise baseline's
last overlap is 56.445 s). The controller avoids again and confirms recovery
at 66 s. The campaign deliberately observes through that later threat.

\section{Raw seeds, limitations and validation}
All fourteen raw seed ODE45 replays have positive sampled clearance. Only
thirteen raw endpoints lie inside the intrinsic terminal core; optimization
still performs local repair where needed. The seed alone is not a safety
witness. Compared with the preceding Gaussian audit, lateral displacement is
smaller in 6/14 cases, and its median decreases from 15.137 to 13.424 m; the
largest displacement increases from 19.983 to 20.365 m. Thus this is not a
uniform improvement in geometric path compactness. The measured contribution
is motion-aware initialization supporting the requested one-model architecture
and successful closed-loop recovery on this family.

An intermediate fourteen-case campaign recovered in thirteen cases. The
15-m/s accelerating head-on seed slightly violated a collision row, inducing a
nonzero PCBF correction. Interpolating from that primary point could not drive
its nonlinear model error to zero: damping shrinks the CLF displacement, not
the primary displacement from the anchor. Adjusting the coupled speed/turn
seed removed that failure without another SCP linearization. Earlier raw-seed
prototypes are exploratory attempts, not additional successful campaigns. The
legacy stress fixture with a tenfold input correction radius also fails from
the new seed. Its regression now explicitly expects a reported failure, not
an inaccurate command or another same-seed linearization. Thus success of the
fourteen default cases is not a universal feasibility claim.

The final focused MATLAB regression contains 388 passing tests in 18 suites, covering target prediction, field frame/time consistency, shifted inputs, fresh initialization on failure, inherited budgets, damping, the single CLF, terminal sets, tire/road-load models and nominal recovery. A persistent solver entry point also fixes test-fixture path restoration invalidating a cached availability flag.

Reproduction uses \texttt{runTimedFlowCampaign(root,out,1500)} and
\texttt{auditFlowInitialization(out)}. Machine-readable results, test outcomes,
source hashes and exact commands are adjacent to this report. Raw traces and
identified archive export copies retain the underlying evidence. The historical
Gaussian audit must be reproduced from its recorded commit; its retired
algorithm is not retained as a controller option.
\end{document}
